\section{Compact monodromy and the descent maps}\label{sec:descent}\label{har:sec:descent}
We prove Corollary~\ref{thm:ordinary} by applying invariant
nonvanishing to the compact factor in the universal cover and then
taking a finite norm. We use the structure statement of \cite[Theorem \textnormal{(Finite cover with compact torus monodromy)}]{CompanionH}, not that companion's nonvanishing conclusion. For a smooth connected projective $X$ with smoothly semipositive $-K_X$, it gives a holomorphic isometric universal-cover splitting
\[
 \widetilde X=\C^c\times V\times F,
\]
where $V$ is a product of compact simply connected Ricci-flat factors with invariant parallel canonical frames and $F$ is a product of compact projective factors with no positive-degree holomorphic forms. After finite index, the deck subgroup acts faithfully by a translation lattice on $\C^c$, trivially on $V$, and through a compact real torus $H$ on $F$. Empty products are points. The associated finite cover $X'$ is smooth projective and need not be a product. These are the only geometric conclusions from the companion that we need.

The compact factor $F$ has $H^j(F,\mathcal O_F)=0$ for $j>0$ by
Hodge symmetry, and its anticanonical metric is smoothly semipositive.
We first descend an invariant section. Afterwards we construct two
additional descent maps, preserving a cohomology class or a twisted
form, to which the ordinary conversions apply on $X'$.

\subsection{Invariant sections and the finite norm}\label{har:subsec:sectiondescent}
\begin{proof}[Proof of Corollary~\ref{thm:ordinary}]
On $F$, the compact torus preserves a very ample line because it preserves its Chern class and $\operatorname{Pic}^0(F)=0$. Its projective action has algebraic-torus Zariski closure: a compact connected abelian subgroup of a projective linear group is conjugate into a diagonal compact torus. Passing through the finite special-linear cover, if necessary, provides a linearization of a very ample power. Invariants of the natural anticanonical line for the compact group and its Zariski closure agree. Corollary~\ref{thm:main} therefore gives a nonzero invariant section of $-mK_F$ for some $m>0$.

Tensor its pullback to $\widetilde X$ with the $m$-th powers of the inverse canonical frames on $\C^c$ and $V$. This nonzero product is invariant under the selected deck subgroup and descends to a section on $X'$. If $F$ is a point, take its section to be $1$ and $m=1$.

The finite \'etale norm \cite[Lemma \textnormal{(Finite \'etale norm)}]{CompanionH} sends this section to a nonzero section of a positive power of $-K_X$. Explicitly, pull to a connected finite Galois refinement, multiply all Galois translates, and descend their invariant product. The product is nonzero on the integral cover, and canonical bundles pull back under \'etale morphisms. Projectivity and GAGA make the resulting holomorphic section algebraic.
\end{proof}

The exponent can increase when common powers are taken on compact factors and when the finite norm is formed. None of these arguments supplies a uniform exponent.

\subsection{Descending a nonzero cohomology class}\label{har:subsec:cohomdescent}

We can also descend cohomology before producing any section.
The invariant Euler characteristic of $F$ is one. Proposition~\ref{har:prop:index} gives, for unbounded \(m\),
nonzero classes in \(H^q(F,-mK_F)^H\), with one fixed \(q\) after
subsequence extraction. Choose a smooth closed Dolbeault representative
and average it over \(H\). Pull it to the product universal cover and
tensor with powers of the constant anticanonical volume on
\(\mathbb C^c\) and the invariant trivializations on \(V\). It is
deck-invariant for the selected subgroup and descends to the finite
cover \(X'\).

The descended class is nonzero. The faithful translation action makes
a compact slice \(F\), at fixed Euclidean and \(V\) coordinates,
embed into \(X'\); no nonidentity lattice element preserves that
slice. Restricting the descended coefficient-valued Dolbeault class
to it, with the product coefficient frames fixed, recovers the chosen
nonzero class on \(F\). Hence \(H^q(X',-mK_{X'})\ne0\) for unbounded
\(m\). Theorem~\ref{thm:alteration-conversion} now produces a positive
anticanonical section on \(X'\). The finite norm of
Section~\ref{har:subsec:sectiondescent} gives a positive anticanonical
section on $X$. Here the invariant cohomology is descended before the
ordinary alteration conversion is applied on $X'$.

\subsection{Descending twisted forms before Iitaka conversion}
\label{har:subsec:formdescent}

There is a direct twisted-form route to the ordinary conversion of
Section~\ref{inj:injectivity}.  Keep the product
$\widetilde X=\mathbb C^c\times V\times F$ and the finite-index
deck subgroup just constructed, whose action is by translations on
$\mathbb C^c$, trivial on $V$, and contained in the compact torus
$H$ on $F$.  Put $P=-K_F$ and $d=\dim F$.  The invariant index has
constant term one.  Hence for a fixed $q$ and unbounded positive
integers $m$ there are nonzero invariant classes in
\[
 H^q(F,mP)^H=H^q(F,K_F+(m+1)P)^H.
\]
Choose an $H$-invariant K\"ahler form.  Smooth semipositive hard
Lefschetz \cite{DPS2001} is an equivariant surjection from
$H^0(F,\Omega_F^{d-q}\otimes (m+1)P)$ onto the displayed group.
Averaging a preimage over $H$ keeps its nonzero image, so supplies
\[
 0\ne\alpha_m\in
 H^0(F,\Omega_F^{d-q}\otimes (m+1)P)^H.
\]
Set $p=d-q$.  Let $\tau_0$ be the translation-invariant inverse
canonical frame on $\mathbb C^c$, and let $\tau_V$ be a product of
the parallel inverse canonical frames on $V$.  Projection to $F$
and the natural inclusion of its pulled-back cotangent bundle into
the product cotangent bundle give the holomorphic twisted form
\[
 \widetilde\alpha_m=
 \operatorname{pr}_F^*\alpha_m\otimes
 (\tau_0\otimes\tau_V)^{m+1}
 \in H^0\bigl(\widetilde X,
       \Omega_{\widetilde X}^{p}\otimes(-K_{\widetilde X})^{m+1}\bigr).
\]
It is nonzero: evaluate $\alpha_m$ on a tangent $p$-tuple at a point
where it is nonzero and use the same tuple in the product; the other
coefficient frames never vanish.  The $m+1$ coefficient exponent
matches the full anticanonical line of the product.

Every factor in this expression is invariant under the selected
deck subgroup.  The resulting form therefore descends through the
covering map to a holomorphic section
\[
 0\ne\beta_m\in
 H^0\bigl(X',\Omega_{X'}^{p}\otimes(-K_{X'})^{m+1}\bigr)
\]
on the associated smooth projective finite \'etale cover $X'$ of
$X$.  Descent follows on evenly covered coordinate neighborhoods,
where invariant local forms glue; its pullback is
$\widetilde\alpha_m$, so the descended form cannot be zero.
For every fixed $r>0$ the same construction is linear and gives
the injection
\[
 H^0(F,\Omega_F^p\otimes rP)^H
 \lhook\joinrel\longrightarrow
 H^0(X',\Omega_{X'}^p\otimes(-rK_{X'})).
\]
Projectivity and GAGA make $\beta_m$ algebraic.  If $F$ is a point,
take $p=q=0$ and $\alpha_m=1$; the same construction applies.

Now apply Corollary~\ref{inj:forms-input} on $X'$ itself to this
fixed differential degree and the unbounded positive twists $m+1$.
It gives a nonzero positive anticanonical section on $X'$.  Pull back
to a finite connected Galois refinement and multiply its translates,
as in Section~\ref{har:subsec:sectiondescent}, to obtain a nonzero
positive anticanonical section on $X$.  Thus this route descends the
twisted forms before using the ordinary Iitaka and injectivity
conversion; the cohomology and alteration route of the preceding
subsection remains a separate construction.

\subsection{The stronger tensor-holonomy alternative}
The Bochner--holonomy argument underlies the structure results of
\cite[Sections~2--3]{DPS1996} and
\cite[proof of Theorem~1.4]{CDP2015}. We include the following tensor
form of that argument as a second input for the same descent. It
excludes all positive covariant tensor powers on an irreducible
non-Ricci-flat factor, which is stronger than excluding holomorphic
differential forms. The proof uses the determinant character of the
irreducible holonomy representation.

\begin{proposition}[The covariant-tensor holonomy dichotomy]
\label{prop:tensor-holonomy}
Let $P$ be a compact simply connected nonflat K\"ahler factor,
irreducible in the Riemannian de Rham sense, with nonnegative Ricci
curvature.  Either $K_P$ has a nonzero
parallel holomorphic volume form, or
\[
 H^0\bigl(P,(\Omega_P^1)^{\otimes m}\bigr)=0
 \qquad\text{for every integer }m>0.
\]
\end{proposition}

\begin{proof}
For a holomorphic covariant tensor, the contracted curvature on its
bundle is the sum, over its tensor slots, of the negative Ricci
operator.  It is nonpositive.  Integration of the Bochner identity
on compact $P$ therefore makes the tensor parallel.  Its value at
a point must consequently be fixed by the holonomy representation.

Let $H$ be the closure of the holonomy group on the complex tangent
space at that point.  It is a compact subgroup of the unitary group.
Simple connectivity makes full holonomy equal to restricted holonomy,
so $H$ is connected.  De Rham irreducibility makes its complex tangent
representation irreducible: a proper invariant complex subspace
would give a proper real parallel subspace and its orthogonal
complement.  Consider the determinant character of this
representation.  If it is trivial, the induced holonomy on the
canonical line is trivial.  Parallel transport of a nonzero canonical
vector gives a global parallel volume form, which is holomorphic.

If the determinant character is nontrivial, connectedness implies
that its differential is nonzero.  This differential is the complex
trace on the Lie algebra $\mathfrak h$.  A compact Lie algebra is the
sum of its center and its derived algebra, and trace vanishes on the
latter.  Hence there is a central $Z\in\mathfrak h$ with nonzero
trace.  By complex irreducibility and Schur's lemma it acts on the
tangent space as a scalar $ia\,\mathrm{Id}$, with $a\in\mathbb R$
nonzero.  Its action on the $m$-fold covariant tensor power is the
nonzero scalar $-mia$.  That tensor power has no $H$-fixed vector.
The parallelism proved in the first paragraph now excludes every
nonzero holomorphic tensor of positive degree.
\end{proof}

This dichotomy supplies a second factor argument, stronger than
vanishing of differential forms.  On its trivial-canonical factors,
a connected group of holomorphic isometries fixes a holomorphic
volume form: it acts trivially on cohomology by isotopy, and the
holomorphic top-form space injects into cohomology by Hodge theory.
On the other factors all positive-degree forms vanish as direct
summands of covariant tensor powers.  Hodge symmetry then gives
vanishing of higher structure-sheaf cohomology.  The compact-monodromy
algebraization, invariant section, and norm construction above apply
to this dichotomy as well.
